\documentclass{article}
\usepackage{graphicx} % Required for inserting images

\usepackage{graphicx}
\usepackage{amsmath, amssymb, amsfonts}
\usepackage{mathtools}
\usepackage{bm}
\usepackage{enumitem}

\usepackage{geometry}
\geometry{margin=1in}

\usepackage{parskip}
\usepackage{hyperref}


\title{Formal Mathematical Phenolgy Metrics}
\author{Parth Thakur}
\date{March 2026}

\begin{document}

\maketitle

\section{Phenology Metrics and Deciduousness Scoring (Formal Definitions)}
\label{sec:phenology-formal}


\subsection{Notation and Data Model}

Let $\mathcal{T}$ denote the set of tracked trees (consensus crowns) and let $\mathcal{O}=\{o_1,\dots,o_T\}$
be the ordered set of orthomosaic indices (OM IDs) used in the analysis.

For each tree $t\in\mathcal{T}$ and orthomosaic $o\in\mathcal{O}$, the pipeline extracts an RGB image patch
(crown chip)
\begin{equation}
P_{t,o}\in\mathbb{R}^{H\times W\times 3}
\end{equation}
by masking the orthomosaic with the consensus crown polygon and cropping to the polygon bounds.
Pixels are expected to be in 8-bit RGB intensity units ($[0,255]$) prior to normalization.

Index pixels by $p\in\{1,\dots,HW\}$. Let $(R_p,G_p,B_p)$ denote the raw (un-normalized) RGB channels for pixel $p$.
After clipping to valid 8-bit range, the implementation scales channels to $[0,1]$:
\begin{equation}
(r_p,g_p,b_p) \;=\; \frac{1}{255}\,\bigl(\mathrm{clip}(R_p,0,255),\,\mathrm{clip}(G_p,0,255),\,\mathrm{clip}(B_p,0,255)\bigr).
\end{equation}

\subsection{Pixel Validity, Masks, and Basic Counts}

\paragraph{Finite and non-zero mask.}
A pixel is considered \emph{valid} if (i) all RGB channel values are finite and (ii) the raw (pre-normalization) RGB sum is strictly positive:
\begin{equation}
\mathbb{1}_{\mathrm{finite}}(p) = \mathbb{1}[R_p, G_p, B_p\ \text{all finite}],\qquad
\mathbb{1}_{\mathrm{nonzero}}(p) = \mathbb{1}[R_p+G_p+B_p>0],
\end{equation}
\begin{equation}
\mathbb{1}_{\mathrm{valid}}(p)=\mathbb{1}_{\mathrm{finite}}(p)\,\mathbb{1}_{\mathrm{nonzero}}(p).
\end{equation}
Define
\begin{equation}
N = HW,\qquad N_{\mathrm{valid}}=\sum_{p=1}^{N} \mathbb{1}_{\mathrm{valid}}(p).
\end{equation}
The valid-pixel fraction is
\begin{equation}
\mathrm{valid\_pixel\_fraction}(t,o)=\frac{N_{\mathrm{valid}}}{N}.
\end{equation}

\paragraph{Small-sample early exit.}
If $N=0$ or $N_{\mathrm{valid}}<20$, feature extraction returns only \texttt{valid\_pixel\_fraction} and no other metrics.

\subsection{Per-pixel Chromatic Signals (GCC, RCC, ExG)}

Define the per-pixel channel sum
\begin{equation}
S_p = r_p+g_p+b_p.
\end{equation}
The implementation uses a numerically safe division operator
$\mathrm{safe\_div}(a,b)=a/(b+\varepsilon)$ with $\varepsilon=10^{-8}$.

\paragraph{Green chromatic coordinate (GCC).}
\begin{equation}
\mathrm{GCC}_p = \frac{g_p}{S_p+\varepsilon}.
\end{equation}
The patch-level mean GCC is the mean over valid pixels:
\begin{equation}
\mathrm{gcc\_mean}(t,o) = \frac{1}{N_{\mathrm{valid}}}\sum_{p=1}^{N} \mathbb{1}_{\mathrm{valid}}(p)\,\mathrm{GCC}_p.
\end{equation}

\paragraph{Red chromatic coordinate (RCC).}
\begin{equation}
\mathrm{RCC}_p = \frac{r_p}{S_p+\varepsilon},
\qquad
\mathrm{rcc\_mean}(t,o) = \frac{1}{N_{\mathrm{valid}}}\sum_{p=1}^{N} \mathbb{1}_{\mathrm{valid}}(p)\,\mathrm{RCC}_p.
\end{equation}

\paragraph{Excess Green (ExG).}
As implemented (no normalization by $S_p$),
\begin{equation}
\mathrm{ExG}_p = 2g_p - r_p - b_p,
\qquad
\mathrm{exg\_mean}(t,o)=\frac{1}{N_{\mathrm{valid}}}\sum_{p=1}^{N}\mathbb{1}_{\mathrm{valid}}(p)\,\mathrm{ExG}_p.
\end{equation}

\subsection{HSV-derived Vegetation and Shadow Fractions}

Let $(h_p,s_p,v_p)\in[0,1]^3$ be the HSV representation of pixel $p$ computed from $(r_p,g_p,b_p)$ \\
using \texttt{matplotlib.colors.rgb\_to\_hsv}.

\paragraph{Shadow (dark) mask.}
A pixel is labeled \emph{dark} if
\begin{equation}
\mathbb{1}_{\mathrm{dark}}(p) = \mathbb{1}[v_p < 0.12].
\end{equation}
The shadow fraction is computed over valid pixels:
\begin{equation}
\mathrm{shadow\_fraction}(t,o)=\frac{1}{N_{\mathrm{valid}}}\sum_{p=1}^{N} \mathbb{1}_{\mathrm{valid}}(p)\,\mathbb{1}_{\mathrm{dark}}(p).
\end{equation}

\paragraph{Vegetation-like mask.}
A pixel is labeled vegetation-like if it is valid and satisfies fixed HSV thresholds:
\begin{equation}
\mathbb{1}_{\mathrm{veg}}(p)=\mathbb{1}_{\mathrm{valid}}(p)\,\mathbb{1}[0.18\le h_p\le 0.48]\,\mathbb{1}[s_p\ge 0.15]\,\mathbb{1}[v_p\ge 0.12].
\end{equation}
The vegetation fraction is
\begin{equation}
\mathrm{veg\_fraction\_hsv}(t,o)=\frac{1}{N_{\mathrm{valid}}}\sum_{p=1}^{N}\mathbb{1}_{\mathrm{veg}}(p).
\end{equation}

\subsection{Channel Statistics and Coefficient of Variation}

All per-channel statistics are computed on the valid pixel set.
Define the valid-pixel index set $\mathcal{V}=\{p:\mathbb{1}_{\mathrm{valid}}(p)=1\}$.

\paragraph{Means.}
\begin{equation}
\mathrm{r\_mean}(t,o)=\frac{1}{|\mathcal{V}|}\sum_{p\in\mathcal{V}} r_p,\quad
\mathrm{g\_mean}(t,o)=\frac{1}{|\mathcal{V}|}\sum_{p\in\mathcal{V}} g_p,\quad
\mathrm{b\_mean}(t,o)=\frac{1}{|\mathcal{V}|}\sum_{p\in\mathcal{V}} b_p.
\end{equation}

\paragraph{Standard deviations.}
Using the population standard deviation (NumPy default, \texttt{ddof=0}):
\begin{equation}
\mathrm{r\_std}(t,o)=\sqrt{\frac{1}{|\mathcal{V}|}\sum_{p\in\mathcal{V}}(r_p-\mathrm{r\_mean})^2},\quad
\mathrm{g\_std}(t,o)=\sqrt{\frac{1}{|\mathcal{V}|}\sum_{p\in\mathcal{V}}(g_p-\mathrm{g\_mean})^2},\quad
\mathrm{b\_std}(t,o)=\sqrt{\frac{1}{|\mathcal{V}|}\sum_{p\in\mathcal{V}}(b_p-\mathrm{b\_mean})^2}.
\end{equation}

\paragraph{Coefficients of variation (CV).}
\begin{equation}
\mathrm{r\_cv}(t,o)=\frac{\mathrm{r\_std}(t,o)}{\mathrm{r\_mean}(t,o)+\varepsilon},\quad
\mathrm{g\_cv}(t,o)=\frac{\mathrm{g\_std}(t,o)}{\mathrm{g\_mean}(t,o)+\varepsilon},\quad
\mathrm{b\_cv}(t,o)=\frac{\mathrm{b\_std}(t,o)}{\mathrm{b\_mean}(t,o)+\varepsilon}.
\end{equation}

\subsection{Grayscale Texture Metrics (Std, Smoothed Std, Entropy, Laplacian Variance)}

\paragraph{Grayscale conversion.}
The pipeline uses a fixed luminance-weighted grayscale:
\begin{equation}
\mathrm{gray}_p = 0.299\,r_p + 0.587\,g_p + 0.114\,b_p.
\end{equation}
A smoothed grayscale image is computed by Gaussian blurring with kernel size $(5,5)$ and standard deviation $0$ (OpenCV default):
\begin{equation}
\mathrm{gray\_smooth} = \mathrm{GaussianBlur}(\mathrm{gray};\,\text{kernel }(5,5)).
\end{equation}

\paragraph{Grayscale standard deviation (raw and smoothed).}
Computed on valid pixels:
\begin{equation}
\mathrm{gray\_std}(t,o)=\mathrm{StdDev}\bigl(\{\mathrm{gray}_p: p\in\mathcal{V}\}\bigr),\qquad
\mathrm{gray\_std\_smooth}(t,o)=\mathrm{StdDev}\bigl(\{\mathrm{gray\_smooth}_p: p\in\mathcal{V}\}\bigr).
\end{equation}

\paragraph{Grayscale entropy.}
The feature \texttt{gray\_entropy} is computed from the valid-pixel grayscale values using a histogram over $[0,1]$ with $B=64$ bins.

Two definitions appear in the notebook:
\begin{itemize}[leftmargin=*, itemsep=2pt]
\item \textbf{Earlier (density) version (overwritten later):}
$h_b$ is the histogram with \texttt{density=True} (bin density). Entropy is
$H=-\sum_b h_b\log_2 h_b$ over bins with $h_b>0$.
\item \textbf{Final (probability) version (the intended definition):}
Let $c_b$ be the bin counts (\texttt{density=False}), $C=\sum_b c_b$, and $p_b=c_b/C$.
The Shannon entropy is
\begin{equation}
\mathrm{gray\_entropy}(t,o) = -\sum_{b=1}^{B} p_b\,\log_2 p_b,\qquad p_b>0.
\end{equation}
\end{itemize}
\noindent
Unless feature extraction is re-run after the entropy-fix cell, exported values may reflect the earlier density-based definition.

\paragraph{Laplacian variance (sharpness).}
The grayscale image is converted to 8-bit
$\mathrm{gray\_u8} = \mathrm{uint8}(\mathrm{clip}(255\,\mathrm{gray},0,255))$.
Then
\begin{equation}
\mathrm{laplacian\_var}(t,o) = \mathrm{Var}\bigl(\Delta\,\mathrm{gray\_u8}\bigr)
\end{equation}
where $\Delta$ is the OpenCV Laplacian operator (second spatial derivative) computed as
\texttt{cv2.Laplacian(gray\_u8, CV\_64F)} and variance is taken over all pixels.

\subsection{GLCM Texture Metrics (Optional)}

If scikit-image texture support is available, the pipeline computes Gray-Level Co-occurrence Matrix (GLCM) properties.

\paragraph{Quantization.}
Pixels are quantized into $L=32$ gray levels:
\begin{equation}
q_p = \mathrm{clip}\Bigl(\bigl\lfloor 31\,\mathrm{gray}_p\bigr\rfloor,\,0,\,31\Bigr)\in\{0,\dots,31\}.
\end{equation}
(Implementation detail: \texttt{(gray * 31).astype(uint8)} with clipping.)

\paragraph{Co-occurrence matrices.}
For distance $d=1$ and angles
\begin{equation}
\Theta = \Bigl\{0,\frac{\pi}{4},\frac{\pi}{2},\frac{3\pi}{4}\Bigr\},
\end{equation}
the GLCM $G^{(\theta)}\in\mathbb{R}^{L\times L}$ is computed with
\texttt{symmetric=True} and \texttt{normed=True}.

\paragraph{GLCM properties.}
Let $\phi_{\mathrm{contrast}}(G)$, $\phi_{\mathrm{homogeneity}}(G)$, and $\phi_{\mathrm{energy}}(G)$
denote \texttt{graycoprops} outputs. The reported metrics are the mean over angles:
\begin{align}
\mathrm{glcm\_contrast}(t,o) &= \frac{1}{|\Theta|}\sum_{\theta\in\Theta} \phi_{\mathrm{contrast}}\bigl(G^{(\theta)}\bigr),\\
\mathrm{glcm\_homogeneity}(t,o) &= \frac{1}{|\Theta|}\sum_{\theta\in\Theta} \phi_{\mathrm{homogeneity}}\bigl(G^{(\theta)}\bigr),\\
\mathrm{glcm\_energy}(t,o) &= \frac{1}{|\Theta|}\sum_{\theta\in\Theta} \phi_{\mathrm{energy}}\bigl(G^{(\theta)}\bigr).
\end{align}
If texture computation fails or scikit-image is unavailable, these values are set to NaN.

\subsection{Observation-level Quality Control (Bad Observation Flag)}

Each $(t,o)$ observation is marked as a bad observation if any of the following holds:
\begin{equation}
\mathrm{is\_bad\_observation}(t,o)=\mathbb{1}\bigl[\mathrm{valid\_pixel\_fraction}(t,o) < 0.60\bigr]
\;\lor\; 
\mathbb{1}\bigl[\mathrm{shadow\_fraction}(t,o) > 0.55\bigr]
\;\lor\;
\mathbb{1}\bigl[\mathrm{laplacian\_var}(t,o) < 25\bigr].
\end{equation}

\subsection{Date-wise Robust Normalization (Illumination Reduction)}

Let $x_{t,o}$ denote any scalar feature (e.g., \texttt{gcc\_mean}) for tree $t$ at orthomosaic $o$.
For each $o$, define the median across trees
\begin{equation}
\mathrm{med}_o = \mathrm{median}_{t\in\mathcal{T}}\, x_{t,o}
\end{equation}
and the median absolute deviation (MAD)
\begin{equation}
\mathrm{MAD}_o = \mathrm{median}_{t\in\mathcal{T}}\,\bigl|x_{t,o}-\mathrm{med}_o\bigr|.
\end{equation}
The pipeline uses a robust z-score with the Gaussian-consistent scaling factor $1.4826$:
\begin{equation}
\mathrm{rz\_date}_{t,o} = \frac{x_{t,o}-\mathrm{med}_o}{1.4826\,\mathrm{MAD}_o + 10^{-8}}.
\end{equation}
Additionally, the date-wise percentile rank is computed (Pandas \texttt{rank(pct=True)} within each $o$):
\begin{equation}
\mathrm{pct\_date}_{t,o} \in (0,1].
\end{equation}

\subsection{Per-tree Temporal Descriptors (Amplitude and Slope)}

Per-tree descriptors are computed on \emph{clean} observations only, i.e.,
\begin{equation}
\mathcal{O}^{\mathrm{clean}}_t = \{o\in\mathcal{O}: \mathrm{is\_bad\_observation}(t,o)=0\}.
\end{equation}
For each metric $x$ (e.g., \texttt{veg\_fraction\_hsv}), the amplitude is
\begin{equation}
\mathrm{amplitude}_t(x) = \max_{o\in\mathcal{O}^{\mathrm{clean}}_t} x_{t,o}\; -\; \min_{o\in\mathcal{O}^{\mathrm{clean}}_t} x_{t,o}.
\end{equation}
If fewer than two clean observations exist, amplitude is recorded as NaN.

For selected metrics, the slope per OM is the least-squares slope of a line fit to $(o, x_{t,o})$ over clean observations,
computed via a first-order polynomial fit:
\begin{equation}
\mathrm{slope\_per\_om}_t(x) = \arg\min_{a,b}\sum_{o\in\mathcal{O}^{\mathrm{clean}}_t} (x_{t,o} - (a\,o + b))^2,
\end{equation}
where the recorded value is $a$.

\subsection{Time-series Construction for Clustering and Leaf-shed Detection}

For a chosen set of metrics $\mathcal{M}$, construct a tree-by-time matrix for each $m\in\mathcal{M}$.
For a fixed tree $t$ and metric $m$, let
\begin{equation}
\tilde{x}^{(m)}_{t} = \bigl( x^{(m)}_{t}(o_{1}),\dots,x^{(m)}_{t}(o_{T})\bigr)\in\mathbb{R}^T
\end{equation}
where bad observations are set to NaN before further processing.

\subsubsection{Linear interpolation with edge fill}

Define $\xi_k=o_k$ and $y_k=\tilde{x}^{(m)}_{t}[k]$.
Let $\mathcal{K}=\{k: y_k\ \text{is finite}\}$.

\begin{itemize}[leftmargin=*, itemsep=2pt]
\item If $|\mathcal{K}|\ge 2$, define $\hat{x}^{(m)}_{t}(\xi)$ as the 1D linear interpolant through the finite points
$\{(\xi_k, y_k): k\in\mathcal{K}\}$ with boundary values held constant outside the observed range.
The interpolated series is
$\bigl(\hat{x}^{(m)}_{t}(o_{1}),\dots,\hat{x}^{(m)}_{t}(o_{T})\bigr)$.
\item If $|\mathcal{K}|=1$, all missing entries are filled with the single observed value.
\item If $|\mathcal{K}|=0$, all entries are filled with $0$.
\end{itemize}

\subsubsection{Per-tree min--max normalization (shape normalization)}

For each tree $t$ and metric $m$, let the interpolated series be $\hat{x}^{(m)}_{t}[1..T]$.
Define
\begin{equation}
\min_t^{(m)}=\min_{k\in\{1,\dots,T\}} \hat{x}^{(m)}_{t}[k],
\qquad
\max_t^{(m)}=\max_{k\in\{1,\dots,T\}} \hat{x}^{(m)}_{t}[k],
\qquad
\mathrm{rng}_t^{(m)}=\max_t^{(m)}-\min_t^{(m)}.
\end{equation}
The normalized series is
\begin{equation}
\bar{x}^{(m)}_{t}[k] =
\begin{cases}
\dfrac{\hat{x}^{(m)}_{t}[k]-\min_t^{(m)}}{\mathrm{rng}_t^{(m)}} & \text{if } \mathrm{rng}_t^{(m)} > 10^{-9},\\
0 & \text{otherwise.}
\end{cases}
\end{equation}
This transformation preserves seasonal \emph{shape} while removing absolute level.

\subsection{Scalar Phenometrics Used for Post-hoc Labelling}

For each tree $t$ and metric $m$ (after interpolation), define:
\begin{align}
\mathrm{amp}_t^{(m)} &= \max_{k} \hat{x}^{(m)}_t[k] - \min_{k} \hat{x}^{(m)}_t[k],\\
k^{\mathrm{pk}}_t(m) &= \arg\max_{k}\,\hat{x}^{(m)}_t[k],\\
k^{\mathrm{tr}}_t(m) &= \arg\min_{k}\,\hat{x}^{(m)}_t[k],\\
\mathrm{peak\_om}_t^{(m)} &= o_{k^{\mathrm{pk}}_t(m)},\\
\mathrm{trough\_om}_t^{(m)} &= o_{k^{\mathrm{tr}}_t(m)}.
\end{align}

The implementation also records a slope with respect to the time index $k\in\{0,\dots,T-1\}$ (not OM ID), via
\begin{equation}
\mathrm{slope}_t^{(m)} = \text{first coefficient of } \mathrm{polyfit}(k,\hat{x}^{(m)}_t[k],1).
\end{equation}

\paragraph{Directional changes in the normalized series.}
Let $d_k = \bar{x}^{(m)}_t[k+1]-\bar{x}^{(m)}_t[k]$ for $k=1,\dots,T-1$.
Discard small changes by thresholding $|d_k|>0.05$; let $\mathcal{D}=\{k: |d_k|>0.05\}$.
Let $s_k=\mathrm{sign}(d_k)$ for $k\in\mathcal{D}$.
The directional-change count is
\begin{equation}
\mathrm{dir\_changes}_t^{(m)} = \sum_{j=1}^{|\mathcal{D}|-1} \mathbb{1}[s_{j+1}\ne s_j].
\end{equation}

\subsection{Leaf-shed Classifier: Deciduousness Score and Phenophase Timing}

This stage operates on the interpolated series. Define the primary vegetation signal
$m_\mathrm{veg}=\texttt{veg\_fraction\_hsv}$ and the supporting color signal
$m_\mathrm{gcc}=\texttt{gcc\_mean}$.
Optionally, a texture signal $m_\mathrm{tex}=\texttt{laplacian\_var}$ is used if present.

\subsubsection{Cross-tree amplitude normalizers (90th percentile)}

For each tree $t$, define raw interpolated amplitudes
\begin{equation}
A_t^{\mathrm{veg}} = \max_{k} \hat{x}^{(m_\mathrm{veg})}_t[k] - \min_{k} \hat{x}^{(m_\mathrm{veg})}_t[k],
\qquad
A_t^{\mathrm{gcc}} = \max_{k} \hat{x}^{(m_\mathrm{gcc})}_t[k] - \min_{k} \hat{x}^{(m_\mathrm{gcc})}_t[k].
\end{equation}
Let $A_{90}^{\mathrm{veg}}$ and $A_{90}^{\mathrm{gcc}}$ denote the 90th percentile across trees (NaNs ignored):
\begin{equation}
A_{90}^{\mathrm{veg}}=\mathrm{percentile}_{90}(\{A_t^{\mathrm{veg}}\}_{t\in\mathcal{T}}),
\qquad
A_{90}^{\mathrm{gcc}}=\mathrm{percentile}_{90}(\{A_t^{\mathrm{gcc}}\}_{t\in\mathcal{T}}),
\end{equation}
with a fallback to $1.0$ if the percentile evaluates to $0$.
If texture is used, define $A_{90}^{\mathrm{tex}}$ analogously for $m_\mathrm{tex}$.

\subsubsection{Per-tree score components}

Within each tree loop, the vegetation series is clipped to $[0,1]$ prior to computing minima and amplitudes:
\begin{equation}
\hat{v}_t[k] = \mathrm{clip}(\hat{x}^{(m_\mathrm{veg})}_t[k],0,1).
\end{equation}
Define
\begin{equation}
\mathrm{veg\_min}_t = \min_{k} \hat{v}_t[k],\qquad
\mathrm{veg\_max}_t = \max_{k} \hat{v}_t[k],\qquad
\mathrm{veg\_amp}_t = \mathrm{veg\_max}_t-\mathrm{veg\_min}_t.
\end{equation}
Let $\mathrm{gcc\_amp}_t=A_t^{\mathrm{gcc}}$.

The component scores are:
\begin{align}
 s_{\mathrm{veg\_amp}}(t) &= \min\Bigl(1,\, \frac{\mathrm{veg\_amp}_t}{A_{90}^{\mathrm{veg}}}\Bigr),\\
 s_{\mathrm{gcc\_amp}}(t) &= \min\Bigl(1,\, \frac{\mathrm{gcc\_amp}_t}{A_{90}^{\mathrm{gcc}}}\Bigr),\\
 s_{\mathrm{depth}}(t) &=
\begin{cases}
\dfrac{\tau_{\mathrm{veg}}-\mathrm{veg\_min}_t}{\tau_{\mathrm{veg}}} & \text{if } \mathrm{veg\_min}_t < \tau_{\mathrm{veg}},\\
0 & \text{otherwise,}
\end{cases}
\end{align}
where the fixed threshold is $\tau_{\mathrm{veg}}=\texttt{VEG\_MIN\_THRESH}=0.45$.
If texture is enabled,
\begin{equation}
 s_{\mathrm{tex}}(t)=\min\Bigl(1,\,\frac{A_t^{\mathrm{tex}}}{A_{90}^{\mathrm{tex}}}\Bigr);
\end{equation}
otherwise $s_{\mathrm{tex}}(t)=0$.

\subsubsection{Deciduousness score and decision rule}

If texture is available, the deciduousness score is
\begin{equation}
\mathrm{DS}(t)=0.35\,s_{\mathrm{veg\_amp}}(t)+0.30\,s_{\mathrm{depth}}(t)+0.25\,s_{\mathrm{gcc\_amp}}(t)+0.10\,s_{\mathrm{tex}}(t).
\end{equation}
If texture is not available, the fallback is
\begin{equation}
\mathrm{DS}(t)=0.40\,s_{\mathrm{veg\_amp}}(t)+0.35\,s_{\mathrm{depth}}(t)+0.25\,s_{\mathrm{gcc\_amp}}(t).
\end{equation}
A tree is classified as deciduous iff
\begin{equation}
\mathrm{is\_deciduous}(t)=\mathbb{1}[\mathrm{DS}(t)\ge \tau_{\mathrm{DS}}],\qquad \tau_{\mathrm{DS}}=\texttt{SCORE\_THRESH}=0.70.
\end{equation}

\subsubsection{Phenophase labels per orthomosaic}

Phenophase labels are derived from the per-tree min--max normalized vegetation series
$\bar{x}^{(m_\mathrm{veg})}_t[k]$ with thresholds
\begin{equation}
\tau_{\mathrm{on}}=\texttt{LEAFON\_NORM}=0.65,\qquad
\tau_{\mathrm{off}}=\texttt{LEAFOFF\_NORM}=0.35.
\end{equation}
For each time index $k$:
\begin{equation}
\mathrm{state}_{t}[k] =
\begin{cases}
\texttt{stable} & \text{if } \mathrm{is\_deciduous}(t)=0,\\
\texttt{leaf\_on} & \text{if } \bar{x}^{(m_\mathrm{veg})}_t[k]\ge \tau_{\mathrm{on}},\\
\texttt{leaf\_off} & \text{if } \bar{x}^{(m_\mathrm{veg})}_t[k]\le \tau_{\mathrm{off}},\\
\texttt{transitioning} & \text{otherwise.}
\end{cases}
\end{equation}

\subsubsection{Event timing extraction (OM indices)}

Let
\begin{equation}
 k^{\mathrm{off}}=\arg\min_{k} \,\hat{v}_t[k],\qquad k^{\mathrm{on}}=\arg\max_{k} \,\hat{v}_t[k].
\end{equation}
The full leaf-off OM and peak leaf-on OM are
\begin{equation}
\mathrm{full\_leaf\_off\_om}(t)=o_{k^{\mathrm{off}}},\qquad
\mathrm{leaf\_on\_peak\_om}(t)=o_{k^{\mathrm{on}}}.
\end{equation}

If $t$ is deciduous, the leaf-off start and leaf-on return events are computed from the categorical state sequence:
\begin{itemize}[leftmargin=*, itemsep=2pt]
\item \textbf{Leaf-off start:} scan backward from $k_{\min}$ to find the last index $k$ with $\mathrm{state}_t[k]=\texttt{leaf\_on}$;
then
$\mathrm{leaf\_off\_start\_om}(t)=o_{\min(k+1,\,T)}$. If none exists, this value is left undefined.
\item \textbf{Leaf-on return:} scan forward from $k_{\min}$ to find the first index $k$ with $\mathrm{state}_t[k]=\texttt{leaf\_on}$;
then
$\mathrm{leaf\_on\_return\_om}(t)=o_k$. If none exists, this value is left undefined.
\end{itemize}
For non-deciduous trees, these event OMs are left undefined.


\end{document}
